\documentclass[tikz,border=4pt]{standalone}
\usepackage{ctex}
\usepackage{tikz}
\usetikzlibrary{arrows.meta,positioning}
\begin{document}
\begin{tikzpicture}[
  >=Latex,
  node distance=7mm and 9mm,
  box/.style={draw,rounded corners,align=center,minimum width=3.4cm,minimum height=0.95cm,font=\small},
  start/.style={draw,rounded corners,very thick,align=center,minimum width=3.8cm,minimum height=1cm,font=\small}
]
\node[start] (q) {先判断题面信号};
\node[box,below left=10mm and 30mm of q] (a) {最大・最小＋区间};
\node[box,below=10mm of q] (b) {两个实根 / 相切};
\node[box,below right=10mm and 30mm of q] (c) {根的位置};
\node[box,below=of a] (a2) {对称轴＋定义域};
\node[box,below=of b] (b2) {判别式 $D$};
\node[box,below=of c] (c2) {$D+$轴+$f(p)$};
\node[box,below=18mm of a2] (d) {$|f(x)|=k$};
\node[box,below=18mm of b2] (e) {恒成立};
\node[box,below=18mm of c2] (f) {二变量 / $x^4,x^2$};
\node[box,below=of d] (d2) {绝对值图像＋水平线};
\node[box,below=of e] (e2) {全实数：全局最值/判别式\\有限区间：区间最值};
\node[box,below=of f] (f2) {消元 或 $t=x^2$\\并重写定义域};
\foreach \u/\v in {q/a,q/b,q/c,a/a2,b/b2,c/c2,a2/d,b2/e,c2/f,d/d2,e/e2,f/f2}
  \draw[->] (\u)--(\v);
\end{tikzpicture}
\end{document}
